%! AP Statistics — Unit 3, Lesson 3.4 — Guided Notes KEY
\documentclass[10pt]{article}
\usepackage{apstats-article}
\usepackage{apstats-key}

\begin{document}

\pageheader{Unit 3, Lesson 3.4}{Guided Notes — KEY}
\namedateperiod

\vspace{0.12in}

\begin{objectivebox}
\small By the end of this lesson, I will be able to\dots\\[4pt]
\ans{Define and distinguish the four major types of bias in sampling.}\\[1pt]
\ans{Explain the mechanism and direction of each bias type.}\\[1pt]
\ans{Identify bias in a real survey scenario and describe how it distorts the estimate.}
\end{objectivebox}

\vspace{0.1in}

\begin{vocabbox}
\termblanklong{Bias} \ans{A systematic tendency for certain responses to be favored over others; causes the estimate to be consistently too high or too low.}
\termblanklong{Voluntary response bias} \ans{Occurs when only self-selected volunteers respond; typically those with extreme opinions.}
\termblanklong{Undercoverage bias} \ans{Part of the population has no chance or a reduced chance of being included in the sample.}
\termblanklong{Nonresponse bias} \ans{Selected individuals who don't respond differ systematically from those who do.}
\termblanklong{Response bias} \ans{Question wording, interviewer effects, or the context of the survey distorts answers.}
\end{vocabbox}

\vspace{0.1in}

\begin{notesbox}{Types of Bias in Sampling}
\small
\textbf{Bias} occurs when certain responses are systematically \ans{favored (over- or under-represented)} over others (DAT-2.E.1).

\vspace{8pt}
\renewcommand{\arraystretch}{2.2}
\begin{tabularx}{\linewidth}{@{} >{\bfseries\color{navy}}p{0.26\linewidth}
    >{\raggedright\arraybackslash}p{0.38\linewidth}
    >{\raggedright\arraybackslash}X @{}}
Voluntary response & Only \ans{self-selected individuals} respond. &
\ans{Call-in radio poll; online vote} \\
Undercoverage & Part of the population has a \ans{lower} chance of being included. &
\ans{Landline-only phone survey misses cell-phone users} \\
Nonresponse & Selected individuals who \ans{refuse} to respond differ from those who do. &
\ans{Mail survey where low-income households don't return forms} \\
Response / question-wording & The \ans{survey} instrument or \ans{social} context distorts answers. &
\ans{"Don't you agree the food is terrible?" leading question} \\
\end{tabularx}
\end{notesbox}

\vspace{0.1in}

\begin{notesbox}{Diagnosing Bias on the AP Exam}
\small
A complete answer must:
\begin{enumerate}[itemsep=4pt]
  \item \textbf{Name} the type: \ans{e.g., "nonresponse bias"}
  \item \textbf{Explain the mechanism}: \ans{e.g., "high-income households are more likely to return the survey, so they are over-represented"}
  \item \textbf{State the direction}: the estimate is too \ans{high/low} because \ans{describe the systematic skew}
\end{enumerate}

\vspace{8pt}
\textbf{Key fact:} A \ans{larger} non-random sample does \textbf{not} fix bias.
More responses of a biased type just give more of the same \ans{distortion}.
\end{notesbox}

\vspace{0.1in}

\begin{practicebox}
\small
\begin{enumerate}[leftmargin=1.5em, itemsep=10pt]
  \item A news channel asks viewers to text ``YES'' or ``NO'' on a political issue.\\
        Bias type: \ans{Voluntary response bias} \quad Direction: \ans{Over-represents people with strong opinions (likely skewed in one direction).}

  \item A political poll only calls landlines; young adults are primarily cell-phone-only.\\
        Bias type: \ans{Undercoverage bias} \quad Direction: \ans{Under-represents younger voters; results skew toward the opinions of older adults.}

  \item A survey asks: ``Would you support cutting wasteful government spending?''\\
        Bias type: \ans{Response bias (question-wording bias)} \quad Direction: \ans{The word "wasteful" is loaded; the estimate of support will be too high.}
\end{enumerate}
\end{practicebox}

\end{document}
